\subsection{Validation using two-stage games from \cite{charness2002understanding}}
\label{sec:validate_prompt}

To validate whether $\bm{\theta}^*$ generalizes to new games, we apply the same prompt template and the values to a new set of more complicated sequential two-stage games from CR---the test set.
Like the validation variants from AR (costless and cycle), these games are plausibly driven by similar underlying mechanisms as the training games, but are still different enough to provide a nontrivial test of generalization.
In the first stage, Person A chooses either a given allocation or lets Person B choose one of two other known allocations.
Person B chooses an allocation but is not informed of Person A's choice---until the payoffs are realized.
For example, in one game, players are shown the following options:
\begin{quote}
Stage 1 (Person A chooses): $\underbrace{(\underbrace{500}_{\mbox{To A}}\:, \:\underbrace{500}_{\mbox{To B}})}_{\mbox{``Left''}}$ $\quad$ vs. $\quad$  $\underbrace{\underbrace{(400\:\:,\:600)\:\text{    vs.  }(700\:\:,\: 300)}_{\mbox{Let Person B choose}}}_{\mbox{``Right''}}$.
\end{quote}
\begin{quote}
Stage 2  (Person B chooses): $\underbrace{(\underbrace{400}_{\mbox{To A}}\:,\: \underbrace{600}_{\mbox{To B}})}_{\mbox{``Left''}}$ $\quad$ vs. $\quad$ $\underbrace{(\underbrace{700}_{\mbox{To A}}\:,\: \underbrace{300}_{\mbox{To B}})}_{\mbox{``Right''}}$.
\end{quote}
Table~\ref{tab:cr_games_results} in Appendix~\ref{app:figs} provides all 20 versions of these two-stage games (each with a different set of payoffs), along with the human results from CR.

As a baseline, we elicit \textsc{Gpt-4o's} responses to these 20 games 150 times each with the temperature set to 1.
We then do the same for the theory-grounded sample $\bm{\theta}^*$---each of the three agents in the mixture plays each game 50 times.

Figure~\ref{fig:cr_test} shows the results. 
The top row shows the responses for the \aisub{s} as Person A, and the bottom row for Person B.
Each column corresponds to a different sample of subjects.
The x-axis shows the setting name and the y-axis shows the mean absolute difference between the fraction of \aisub{s} choosing ``Left'' and the fraction of human subjects choosing ``Left'' in \citeauthor{charness2002understanding}.
The difference between the baseline (red) and selected \aisub{s} (blue) is substantial.
The MAE between the baseline and the human subjects as Player A (\CRTestBaselineGameAMAE) is three times larger than that for the optimized agents relative to the humans (\CRTestOptGameAMAE).
The difference in MAE is twice as large for Player B (\CRTestBaselineGameBMAE{} vs. \CRTestOptGameBMAE).

\begin{figure}[h!]
\begin{center}
\begin{minipage}{\textwidth}
\caption{Distances between human and \aisub{s} for the two-stage dictator games} 
\label{fig:cr_test}
\vspace*{-0.15in}
\centering
\includegraphics[width=.9\textwidth]{plots/cr_bilater_dict.pdf} 
\end{minipage}
\end{center}
\begin{footnotesize}
\begin{singlespace}
\vspace*{-0.05in}
\emph{Notes:}
This figure reports the results of replications of the sequential two-player games from \cite{charness2002understanding} with \aisub{s}.
Each row shows responses from either Person A (left) or Person B (right), while each column corresponds to a different set of subjects. 
The x-axis shows the game, and the y-axis shows the mean absolute difference between the fraction of \aisub{s} choosing ``Left'' and the fraction of human subjects choosing ``Left'' in \citeauthor{charness2002understanding}.
The left column displays the baseline \aisub{s} (red), the middle column is the selected \aisub{s} (blue), and the right column shows the atheoretical \aisub{s} (yellow).  
The horizontal dashed lines show the mean absolute error in each pane.
\end{singlespace}
\end{footnotesize}
\end{figure}

This predictive improvement is robust across settings.
In $\CRTestNOptBetterThanBaseline$ of the 40 total decisions (20 settings, each played as Person A and Person B), the optimized theory-grounded agents more accurately matched human behavior than the baseline AI. 
Never was the absolute error in a game larger than 0.50 for the optimized theory-grounded agents, which is less than the MAE for the baseline AI as player A.
